% GENERATED FILE. Edit canonical lessons, then run npm run generate:books.
% canonical-source-hash: fnv1a64:4987218e9a951f06
% canonical-lessons: RU-C320-volshebnik, RU-C320-vedma, RU-C320-chudovishche, RU-C320-drakon, RU-C320-prevrashchatsya

\chapter{Fairy-tale creatures}
\label{ch:ru-fairy-tale-creatures}
\hlchaptermodality{320}

\begin{chapteropening}
\textbf{By the end of this chapter:} \emph{I can talk about a wizard, a witch, a monster and a dragon, and say that someone turned into something else.}

Complete the last lesson of chapter 320: I can talk about a wizard, a witch, a monster and a dragon, and say that someone turned into something else.
\end{chapteropening}

\section[\texorpdfstring{\ru{волшебник} (volshébnik)}{волшебник (volshébnik)}]{\ru{волшебник} (volshébnik) — a wizard, a magician}

\label{lesson:RU-C320-volshebnik}



\begin{quote}
Before the new one: say the Russian for a storyteller, then the Russian for to make up.
\end{quote}

\subsection*{You'll want to know: \ru{волшебник}}
\textbf{\ru{волшебник}} (\emph{volshébnik}) — \textquotedblleft{}a wizard, a magician\textquotedblright{}.

\begin{quote}
\ru{В} \ru{сказке} \ru{волшебник} \ru{помог} \ru{бедной} \ru{девушке}.
\end{quote}

\begin{quote}
\emph{V skázke volshébnik pomóg bédnoy dévushke.}
\end{quote}

\begin{quote}
In the tale a wizard helped a poor girl.
\end{quote}

\emph{\ru{Волшебный}} — magic, magical.

\begin{cousinweb}
\emph{\ru{волшебник}} is built on \emph{\ru{волшебный}} \textquotedblleft{}magic\textquotedblright{} and \emph{-\ru{ик}}, from \emph{\ru{волшба}} \textquotedblleft{}sorcery\textquotedblright{}, Old East Slavic \emph{vŭlšĭba} \textquotedblleft{}magic, sorcery\textquotedblright{} (EtymDB).
\end{cousinweb}

\begin{grammarlens}[title={What you've built}]
One more word for telling, explaining and describing at B1.
\end{grammarlens}

\subsection*{Your turn}
\begin{itemize}
\raggedright
  \item \emph{Say it:} \emph{volshébnik}
  \item \emph{Say it:} \emph{volshébnik}, once more
  \item \emph{Say it:} a sentence of your own with \emph{volshébnik}
  \item \emph{Recall:} say the Russian for a storyteller, then the Russian for to make up, then say \emph{volshébnik} again
\end{itemize}

\begin{tcolorbox}[breakable,colback=teal!4,colframe=teal!35!black,title={Before you move on}]
What does \ru{волшебник} mean? (\textquotedblleft{}A wizard, a magician\textquotedblright{}.) Say it once more.
\end{tcolorbox}
\section[\texorpdfstring{\ru{ведьма} (védma)}{ведьма (védma)}]{\ru{ведьма} (védma) — a witch}

\label{lesson:RU-C320-vedma}



\begin{quote}
Before the new one: say the Russian for to make up, then the Russian for a wizard.
\end{quote}

\subsection*{You'll want to know: \ru{ведьма}}
\textbf{\ru{ведьма}} (\emph{védma}) — \textquotedblleft{}a witch\textquotedblright{}.

\begin{quote}
\ru{Старая} \ru{ведьма} \ru{жила} \ru{в} \ru{лесу}.
\end{quote}

\begin{quote}
\emph{Stáraya védma zhilá v lesú.}
\end{quote}

\begin{quote}
An old witch lived in the forest.
\end{quote}

Feminine.

\begin{cousinweb}
\emph{\ru{ведьма}} is inherited through Old East Slavic \emph{vědĭma} from Proto-Slavic \emph{vědĭma} (EtymDB), built on the root of \emph{\ru{ведать}} \textquotedblleft{}to know\textquotedblright{}: a woman who knows.
\end{cousinweb}

\begin{grammarlens}[title={What you've built}]
One more word for telling, explaining and describing at B1.
\end{grammarlens}

\subsection*{Your turn}
\begin{itemize}
\raggedright
  \item \emph{Say it:} \emph{védma}
  \item \emph{Say it:} \emph{védma}, once more
  \item \emph{Say it:} a sentence of your own with \emph{védma}
  \item \emph{Recall:} say the Russian for to make up, then the Russian for a wizard, then say \emph{védma} again
\end{itemize}

\begin{tcolorbox}[breakable,colback=teal!4,colframe=teal!35!black,title={Before you move on}]
What does \ru{ведьма} mean? (\textquotedblleft{}A witch\textquotedblright{}.) Say it once more.
\end{tcolorbox}
\section[\texorpdfstring{\ru{чудовище} (chudóvishche)}{чудовище (chudóvishche)}]{\ru{чудовище} (chudóvishche) — a monster}

\label{lesson:RU-C320-chudovishche}



\begin{quote}
Before the new one: say the Russian for a wizard, then the Russian for a witch.
\end{quote}

\subsection*{You'll want to know: \ru{чудовище}}
\textbf{\ru{чудовище}} (\emph{chudóvishche}) — \textquotedblleft{}a monster\textquotedblright{}.

\begin{quote}
\ru{Чудовище} \ru{жило} \ru{в} \ru{глубокой} \ru{пещере}.
\end{quote}

\begin{quote}
\emph{Chudóvishche zhílo v glubókoy peshchére.}
\end{quote}

\begin{quote}
The monster lived in a deep cave.
\end{quote}

Neuter, though it ends in \emph{-\ru{е}}: \emph{\ru{страшное} \ru{чудовище}}.

\begin{cousinweb}
\emph{\ru{чудовище}} is built on \emph{\ru{чудо}} \textquotedblleft{}a wonder, a marvel\textquotedblright{}, an inherited Slavic word, with the ending \emph{-\ru{овище}}; EtymDB lists no parent.
\end{cousinweb}

\begin{grammarlens}[title={What you've built}]
One more word for telling, explaining and describing at B1.
\end{grammarlens}

\subsection*{Your turn}
\begin{itemize}
\raggedright
  \item \emph{Say it:} \emph{chudóvishche}
  \item \emph{Say it:} \emph{chudóvishche}, once more
  \item \emph{Say it:} a sentence of your own with \emph{chudóvishche}
  \item \emph{Recall:} say the Russian for a wizard, then the Russian for a witch, then say \emph{chudóvishche} again
\end{itemize}

\begin{tcolorbox}[breakable,colback=teal!4,colframe=teal!35!black,title={Before you move on}]
What does \ru{чудовище} mean? (\textquotedblleft{}A monster\textquotedblright{}.) Say it once more.
\end{tcolorbox}
\section[\texorpdfstring{\ru{дракон} (drakón)}{дракон (drakón)}]{\ru{дракон} (drakón) — a dragon}

\label{lesson:RU-C320-drakon}



\begin{quote}
Before the new one: say the Russian for a witch, then the Russian for a monster.
\end{quote}

\subsection*{You'll want to know: \ru{дракон}}
\textbf{\ru{дракон}} (\emph{drakón}) — \textquotedblleft{}a dragon\textquotedblright{}.

\begin{quote}
\ru{В} \ru{сказке} \ru{дракон} \ru{охранял} \ru{золото}.
\end{quote}

\begin{quote}
\emph{V skázke drakón okhranyál zóloto.}
\end{quote}

\begin{quote}
In the tale a dragon guarded gold.
\end{quote}

Masculine. A Russian folk dragon is \emph{\ru{змей}}.

\begin{cousinweb}
\emph{\ru{дракон}} goes back to Greek \emph{drákōn} \textquotedblleft{}a serpent, a dragon\textquotedblright{} (EtymDB).
\end{cousinweb}

\begin{grammarlens}[title={What you've built}]
One more word for telling, explaining and describing at B1.
\end{grammarlens}

\subsection*{Your turn}
\begin{itemize}
\raggedright
  \item \emph{Say it:} \emph{drakón}
  \item \emph{Say it:} \emph{drakón}, once more
  \item \emph{Say it:} a sentence of your own with \emph{drakón}
  \item \emph{Recall:} say the Russian for a witch, then the Russian for a monster, then say \emph{drakón} again
\end{itemize}

\begin{tcolorbox}[breakable,colback=teal!4,colframe=teal!35!black,title={Before you move on}]
What does \ru{дракон} mean? (\textquotedblleft{}A dragon\textquotedblright{}.) Say it once more.
\end{tcolorbox}
\section[\texorpdfstring{\ru{превращаться} (prevrashchátsya)}{превращаться (prevrashchátsya)}]{\ru{превращаться} (prevrashchátsya) — to turn into, to become}

\label{lesson:RU-C320-prevrashchatsya}



\begin{quote}
Before the new one: say the Russian for a monster, then the Russian for a dragon.
\end{quote}

\subsection*{You'll want to know: \ru{превращаться}}
\textbf{\ru{превращаться}} (\emph{prevrashchátsya}) — \textquotedblleft{}to turn into, to become\textquotedblright{}.

\begin{quote}
\ru{Ночью} \ru{ведьма} \ru{превращалась} \ru{в} \ru{чёрную} \ru{кошку}.
\end{quote}

\begin{quote}
\emph{Nóchyu védma prevrashchálas v chyórnuyu kóshku.}
\end{quote}

\begin{quote}
At night the witch turned into a black cat.
\end{quote}

Imperfective; perfective \emph{\ru{превратиться}} (\emph{prevratítsya}), \emph{\ru{превращусь}}. It takes \emph{\ru{в}} + accusative.

\begin{cousinweb}
\emph{\ru{превращаться}} is built on \emph{\ru{пре}-} and \emph{\ru{вращать}} \textquotedblleft{}to turn\textquotedblright{}, which EtymDB marks as borrowed from Church Slavonic \emph{vraščati}, from Proto-Slavic \emph{vorťati}; the native cousin is \emph{\ru{воротить}}, Proto-Slavic \emph{vortiti} (EtymDB).
\end{cousinweb}

\begin{grammarlens}[title={What you've built}]
One more word for telling, explaining and describing at B1.
\end{grammarlens}

\subsection*{Your turn}
\begin{itemize}
\raggedright
  \item \emph{Say it:} \emph{prevrashchátsya}
  \item \emph{Say it:} \emph{prevrashchátsya}, once more
  \item \emph{Say it:} a sentence of your own with \emph{prevrashchátsya}
  \item \emph{Recall:} say the Russian for a monster, then the Russian for a dragon, then say \emph{prevrashchátsya} again
\end{itemize}

\begin{tcolorbox}[breakable,colback=teal!4,colframe=teal!35!black,title={Before you move on}]
What does \ru{превращаться} mean? (\textquotedblleft{}To turn into, to become\textquotedblright{}.) Say it once more.
\end{tcolorbox}
